\chapter{La derivada de Malliavin}\label{cap:derivada}

\begin{panorama}
Definimos la derivada $D$ en las variables suaves, probamos la fórmula de integración por
partes fundamental, la cerrabilidad de $D$ y construimos los espacios de Sobolev--Watanabe
$\D^{k,p}$. Establecemos la regla de la cadena, la expresión de $D$ en el caos de Wiener (incluida
la fórmula de Stroock) y calculamos la derivada de funcionales típicos: integrales de
Wiener, el supremo del browniano y esperanzas condicionales.
\end{panorama}

Trabajamos con un proceso isonormal $W$ sobre $\Hf$, y $\F=\sigma(W)$.

\section{Definición en las variables suaves}

\begin{definicion}
Para $F=f(W(h_1),\dots,W(h_n))\in\Sc$, la \emph{derivada de Malliavin} es la variable
con valores en $\Hf$
\[
  DF=\sum_{i=1}^n\partial_if\big(W(h_1),\dots,W(h_n)\big)\,h_i.
\]
\end{definicion}

Es fácil ver que $DF$ no depende de la representación de $F$ (si $F$ admite dos representaciones,
se reducen a una común con $h_i$ linealmente independientes). En el espacio de Wiener clásico,
$\Hf=L^2([0,T])$, $DF$ es un proceso $(D_tF)_{t\in[0,T]}$ y
\[
  \ip{DF}{\dot h}_{L^2}=\int_0^TD_tF\,\dot h(t)\,\dd t=\frac{\dd}{\dd\varepsilon}F(\omega+\varepsilon h)\Big|_{\varepsilon=0},\qquad h=\textstyle\int_0^\cdot\dot h,
\]
pues $W(g)(\omega+\varepsilon h)=W(g)(\omega)+\varepsilon\ip{g}{\dot h}$. Es la \emph{derivada de Fréchet en las direcciones de
Cameron--Martin}.

\begin{ejemplo}
\begin{enumerate}[label=(\alph*)]
  \item $DW(h)=h$. En particular $D_tW_s=\ind_{[0,s]}(t)$.
  \item $D_t(W_s^2)=2W_s\ind_{[0,s]}(t)$; $D_t\int_0^TW_s\,\dd s=T-t$.
  \item $D\,e^{W(h)}=e^{W(h)}h$.
\end{enumerate}
\end{ejemplo}

\begin{intuicion}
$D_tF$ mide cuánto cambia $F$ cuando se perturba el ruido \emph{en el instante $t$}: es la respuesta
de $F$ a un ``impulso'' $\delta_t$ en $\dot W$. Por ejemplo, $W_s$ responde con 1 a impulsos anteriores a $s$
y con 0 a impulsos posteriores. Si $F$ es la solución de una EDE en el instante $T$, $D_tX_T$
es la propagación de una perturbación desde $t$ hasta $T$ (Capítulo~\ref{cap:hormander}).
\end{intuicion}

\section{Integración por partes y cerrabilidad}

\begin{lema}[Integración por partes fundamental]\label{lem:ibp}
Para $F\in\Sc$ y $h\in\Hf$,
\[
  \E\ip{DF}{h}_\Hf=\E\big[F\,W(h)\big].
\]
\end{lema}

\begin{proof}
Por homogeneidad podemos suponer $\norm{h}=1$, y podemos escribir $F=
f(W(e_1),\dots,W(e_n))$ con $\{e_i\}$ ortonormales y $e_1=h$. Entonces $(W(e_1),\dots,W(e_n))\sim N(0,I_n)$
y, por el Lema~\ref{lem:stein},
\[
  \E\ip{DF}{h}=\E[\partial_1f(W(e))]=\int_{\R^n}\partial_1f(x)\phi_n(x)\,\dd x=\int_{\R^n}f(x)x_1\phi_n(x)\,\dd x=\E[FW(e_1)].
  \qedhere
\]
\end{proof}

\begin{lema}\label{lem:ibp2}
Para $F,G\in\Sc$ y $h\in\Hf$, $\E\big[G\ip{DF}{h}\big]=\E\big[-F\ip{DG}{h}+FGW(h)\big]$.
\end{lema}

\begin{proof}
Aplicar el Lema~\ref{lem:ibp} a $FG$ y usar la regla de Leibniz $D(FG)=F\,DG+G\,DF$,
obvia en $\Sc$.
\end{proof}

\begin{proposicion}[Cerrabilidad]\label{prop:cerrabilidad}
Para $p\ge1$, el operador $D:\Sc\subset L^p(\Omega)\to L^p(\Omega;\Hf)$ es
cerrable.
\end{proposicion}

\begin{proof}
Sea $F_n\in\Sc$ con $F_n\to0$ en $L^p$ y $DF_n\to\eta$ en $L^p(\Omega;\Hf)$. Hay que ver $\eta=0$. Sea
$G\in\Sc$ tal que $GW(h)$ esté acotada (por ejemplo $G=g\,e^{-\varepsilon W(h)^2}$ con $g\in\Sc_b$). Por el Lema~\ref{lem:ibp2},
\[
  \E\big[G\ip{\eta}{h}\big]=\lim_n\E\big[G\ip{DF_n}{h}\big]=\lim_n\E\big[-F_n\ip{DG}{h}+F_nGW(h)\big]=0,
\]
porque $\ip{DG}{h}$ y $GW(h)$ están acotadas. Tales $G$ son densas en $L^q$ (tomando $\varepsilon\to0$), luego
$\ip{\eta}{h}=0$ c.s.\ para todo $h$ de un subconjunto numerable denso, y $\eta=0$.
\end{proof}

\section{Espacios de Sobolev--Watanabe}

\begin{definicion}
Para $p\ge1$, $\D^{1,p}$ es el dominio de la clausura de $D$ en $L^p$, es decir, la
compleción de $\Sc$ respecto a
\[
  \norm{F}_{1,p}=\big(\E\abs{F}^p+\E\norm{DF}_\Hf^p\big)^{1/p}.
\]
Iterando, $D^kF$ es una variable con valores en $\Hf^{\otimes k}$, y $\D^{k,p}$ es la compleción con $\norm{F}_{k,p}^p=
\E\abs{F}^p+\sum_{j=1}^k\E\norm{D^jF}^p_{\Hf^{\otimes j}}$. Se definen análogamente $\D^{k,p}(V)$ para $V$-valuadas, con $V$ un
Hilbert, y
\[
  \D^\infty=\bigcap_{k,p}\D^{k,p}.
\]
\end{definicion}

$\D^{1,2}$ es un espacio de Hilbert. Se tiene $\D^{k,p}\subset\D^{k',p'}$ si $k\ge k'$, $p\ge p'$. Los elementos de $\D^\infty$
son los ``funcionales suaves de Wiener''. Veremos que, gracias a las desigualdades de Meyer, $\D^{k,p}$
coincide con el dominio de $(I-L)^{k/2}$ en $L^p$.

\section{Regla de la cadena}

\begin{proposicion}[Regla de la cadena]\label{prop:cadena}
Sea $\varphi\in C^1(\R^m)$ con derivadas parciales acotadas y
$F=(F^1,\dots,F^m)$ con $F^i\in\D^{1,p}$. Entonces $\varphi(F)\in\D^{1,p}$ y
\[
  D\varphi(F)=\sum_{i=1}^m\partial_i\varphi(F)\,DF^i.
\]
Si solo se supone $\varphi$ globalmente Lipschitz y $p>1$, entonces $\varphi(F)\in\D^{1,p}$ y existe un vector aleatorio
$G$ con $\abs{G}\le\norm{\varphi}_{\mathrm{Lip}}$ tal que $D\varphi(F)=\sum_iG_iDF^i$. Si además la ley de $F$ es absolutamente
continua, $G=\nabla\varphi(F)$.
\end{proposicion}

\begin{proof}
Para $F^i\in\Sc$ y $\varphi\in C_b^1$ es la regla de la cadena clásica. En general se aproxima
$F^i$ por $F^i_n\in\Sc$ en $\D^{1,p}$; entonces $\varphi(F_n)\to\varphi(F)$ en $L^p$ y $\sum\partial_i\varphi(F_n)DF^i_n\to\sum\partial_i\varphi(F)DF^i$ en
$L^p(\Hf)$ (convergencia dominada sobre una subsucesión con convergencia c.s.); por cerrabilidad,
se concluye. Para $\varphi$ Lipschitz, se regulariza $\varphi_\varepsilon=\varphi*\rho_\varepsilon$; la sucesión $D\varphi_\varepsilon(F)$ está acotada en
$L^p(\Hf)$, que es reflexivo para $p>1$, y un límite débil identifica la derivada.
\end{proof}

\begin{ejemplo}
Si $F\in\D^{1,2}$, entonces $\abs{F}\in\D^{1,2}$ y $D\abs{F}=\sgn(F)DF$; $F^+\in\D^{1,2}$ y $DF^+=
\ind_{\{F>0\}}DF$.
\end{ejemplo}

\begin{ejemplo}[Una variable no derivable]\label{ej:no-derivable}
$F=\ind_{\{W(h)>0\}}\notin\D^{1,2}$ si $h\ne0$. Si lo fuera, por
localidad (Proposición~\ref{prop:localidad}), $DF=0$ c.s.\ en $\{F=0\}\cup\{F=1\}=\Omega$, y por el Corolario~\ref{cor:constante},
$F$ sería constante. Los indicadores de sucesos no triviales nunca son Malliavin-diferenciables.
Esta es la razón de fondo de por qué las griegas de opciones digitales requieren integración por
partes (Capítulo~\ref{cap:griegas}).
\end{ejemplo}

\section{La derivada en el caos de Wiener}

\begin{teorema}\label{teo:derivada-caos}
Sea $F=\sum_nI_n(f_n)\in L^2(\Omega)$. Entonces $F\in\D^{1,2}$ si y solo si
\[
  \sum_{n=1}^\infty n\,n!\,\norm{f_n}^2_{\Hf^{\otimes n}}<\infty,
\]
y en tal caso, en el caso $\Hf=L^2(T,\mu)$,
\[
  D_tF=\sum_{n=1}^\infty nI_{n-1}\big(f_n(\cdot,t)\big),\qquad\E\norm{DF}^2_\Hf=\sum_{n\ge1}n\,n!\,\norm{f_n}^2.
\]
\end{teorema}

\begin{proof}
Para $F=I_n(h^{\otimes n})=H_n(W(h))$ con $\norm{h}=1$: $DF=H_n'(W(h))h=
nH_{n-1}(W(h))h=nI_{n-1}(h^{\otimes(n-1)})h$, que es la fórmula. Por linealidad y densidad vale para todo
$f_n\in\Hf^{\odot n}$, y $\E\norm{DI_n(f_n)}^2=n^2(n-1)!\norm{f_n}^2=n\,n!\norm{f_n}^2$. Además, $DI_n(f_n)$ y $DI_m(f_m)$ son
ortogonales en $L^2(\Omega;\Hf)$ si $n\ne m$. Si la serie converge, las sumas parciales $F_N=\sum_{n\le N}I_n(f_n)$
convergen en $\D^{1,2}$ y, por cerrabilidad, $F\in\D^{1,2}$. Recíprocamente, si $F\in\D^{1,2}$, se usa que $D$
conmuta con las proyecciones $J_n$ (salvo desplazamiento de índice), lo que se comprueba en $\Sc$, y
la serie es $\E\norm{DF}^2<\infty$.
\end{proof}

\begin{corolario}[Fórmula de Stroock]\label{cor:stroock}
Si $F\in\D^{n,2}$, el núcleo del $n$-ésimo caos es $f_n=
\frac1{n!}\E[D^nF]$.
\end{corolario}

\begin{proof}
Aplicando $n$ veces el Teorema~\ref{teo:derivada-caos}, $D^nF=\sum_{m\ge n}\frac{m!}{(m-n)!}I_{m-n}(f_m)$, cuya esperanza es el término $m=n$: $n!f_n$.
\end{proof}

\begin{ejemplo}
$F=W_1^3$. $\E F=0$, $\E DF=\E[3W_1^2]\ind_{[0,1]}=3\cdot\ind_{[0,1]}$, $\E D^2F=0$, $\E D^3F=6\cdot\ind_{[0,1]^3}$.
Así $W_1^3=I_1(3\cdot\ind)+I_3(\ind^{\otimes3})=3W_1+H_3(W_1)$. Correcto: $H_3(x)=x^3-3x$.
\end{ejemplo}

\section{Propiedades estructurales}

\begin{corolario}\label{cor:constante}
Si $F\in\D^{1,2}$ y $DF=0$, entonces $F=\E F$ es constante.
\end{corolario}

\begin{proof}
$0=\E\norm{DF}^2=\sum n\,n!\norm{f_n}^2$ obliga a $f_n=0$ para $n\ge1$.
\end{proof}

\begin{proposicion}[Localidad]\label{prop:localidad}
Si $F\in\D^{1,1}$, entonces $DF=0$ c.s.\ en el conjunto $\{F=0\}$.
Más en general, si $F=G$ en un suceso $A$, entonces $DF=DG$ c.s.\ en $A$.
\end{proposicion}

\begin{proof}[Esbozo]
Sea $\psi_\varepsilon(x)=\int_0^x\ind_{[-\varepsilon,\varepsilon]}(y)\,\dd y$, Lipschitz con $\psi_\varepsilon'=\ind_{[-\varepsilon,\varepsilon]}$ c.t.p.\ y $\abs{\psi_\varepsilon}\le\varepsilon$. Por la regla
de la cadena (versión Lipschitz), $D\psi_\varepsilon(F)=\ind_{\{\abs{F}\le\varepsilon\}}DF$. Para $u$ suave con valores en $\Hf$, la
dualidad del Capítulo~\ref{cap:divergencia} da
\[
  \Big|\E\big[\ind_{\{\abs{F}\le\varepsilon\}}\ip{DF}{u}\big]\Big|=\abs{\E[\psi_\varepsilon(F)\delta(u)]}\le\varepsilon\,\E\abs{\delta(u)}\xrightarrow[\varepsilon\to0]{}0.
\]
Por convergencia dominada, $\E[\ind_{\{F=0\}}\ip{DF}{u}]=0$ para todo $u$ de un conjunto denso, de donde
$\ind_{\{F=0\}}DF=0$.
\end{proof}

\begin{proposicion}[Derivada de una esperanza condicional]\label{prop:derivada-esp-cond}
Sea $\Hf=L^2(T,\mu)$, $A\subset T$ medible
y $\F_A=\sigma\{W(h):\ \operatorname{supp}h\subset A\}$. Si $F\in\D^{1,2}$, entonces $\cond{F}{\F_A}\in\D^{1,2}$ y
\[
  D_t\cond{F}{\F_A}=\cond{D_tF}{\F_A}\ind_A(t).
\]
En particular, con $A=[0,s]$ en el espacio de Wiener: si $F$ es $\F_s$-medible, $D_tF=0$ para $t>s$.
\end{proposicion}

\begin{proof}
Si $F=I_n(f_n)$, entonces $\cond{F}{\F_A}=I_n(f_n\ind_A^{\otimes n})$ (por Proposición~\ref{prop:hermite}(iii), basta
comprobarlo para $f_n=h^{\otimes n}$, donde se descompone $h=h\ind_A+h\ind_{A^c}$, cuyas integrales son
independientes). Aplicando el Teorema~\ref{teo:derivada-caos},
\[
  D_t\cond{F}{\F_A}=nI_{n-1}\big(f_n(\cdot,t)\ind_A^{\otimes(n-1)}\big)\ind_A(t)=\cond{D_tF}{\F_A}\ind_A(t).
\]
Se concluye sumando sobre $n$.
\end{proof}

\section{Un ejemplo no trivial: el supremo del browniano}

\begin{proposicion}\label{prop:supremo}
Sea $M=\sup_{t\in[0,1]}W_t$. Entonces $M\in\D^{1,2}$ y $D_sM=\ind_{[0,\tau]}(s)$, donde $\tau$ es
el único instante \textup{(c.s.)} en que se alcanza el máximo.
\end{proposicion}

\begin{proof}
Sea $\{t_i\}$ denso en $[0,1]$ y $M_n=\max_{i\le n}W_{t_i}$. La aplicación $\varphi_n(x)=\max_{i\le n}x_i$
es Lipschitz, y por la Proposición~\ref{prop:cadena}, $D_sM_n=\ind_{[0,\tau_n]}(s)$ con $\tau_n$ el $t_i$ donde se alcanza el máximo
(único c.s.\ por continuidad absoluta). Así $\norm{DM_n}^2_\Hf=\tau_n\le1$, y $M_n\to M$ en $L^2$. Como $\{M_n\}$
está acotada en $\D^{1,2}$, una subsucesión converge débilmente y, por cerrabilidad, $M\in\D^{1,2}$. La
unicidad del maximizante (consecuencia de la propiedad de Markov) da $\tau_n\to\tau$ c.s., e identifica
$DM=\ind_{[0,\tau]}$.
\end{proof}

\begin{intuicion}
Perturbar el ruido antes de $\tau$ desplaza toda la trayectoria posterior, incluido el máximo;
perturbarlo después de $\tau$ no cambia el máximo (a primer orden). El supremo no es
diferenciable en sentido Fréchet en $C_0$, pero sí en el sentido de Malliavin: la derivada solo
se pide en direcciones de Cameron--Martin y en $L^2$.
\end{intuicion}

\section{Caracterización por trayectorias}
En el espacio de Wiener clásico, $\D^{1,p}$ tiene una descripción intrínseca debida a Kusuoka y
Sugita: $F\in\D^{1,p}$, $p>1$, si y solo si (i) para cada $h\in H$ existe una versión de $F$ tal
que $\varepsilon\mapsto F(\omega+\varepsilon h)$ es absolutamente continua (\emph{ray absolute continuity}); y (ii) existe $DF\in
L^p(\Omega;H)$ tal que $\frac1\varepsilon(F(\omega+\varepsilon h)-F(\omega))\to\ip{DF}{h}$ en probabilidad para todo $h$ (\emph{diferenciabilidad
estocástica de Gâteaux}). En particular, todo funcional $F$ que sea $H$-Lipschitz, es decir, con
$\abs{F(\omega+h)-F(\omega)}\le C\norm{h}_H$, está en $\D^{1,p}$ con $\norm{DF}_H\le C$.
